% Paper A -- ACIIDS 2027 submission draft (Springer LNCS/LNAI), method paper, version 2.
% All numbers come from results.tex, generated by make_results.py from repository files.
\documentclass[runningheads]{llncs}
\usepackage[T1]{fontenc}
\usepackage{graphicx}
\usepackage{amsmath,amssymb}
\usepackage{booktabs}
\usepackage{xcolor}
\usepackage{algorithm}
\usepackage[noend]{algpseudocode}
\usepackage[hidelinks]{hyperref}
\newif\ifpendshort
% \PENDING{key}: placeholder for a value that does not exist yet (RIS-relevant pilot, NSGA-II).
\newcommand{\PENDING}[1]{\textcolor{red}{\textbf{\ifpendshort P\else[#1]\fi}}}
\newif\ifbench
% AUTO-GENERATED by make_results.py -- do not edit by hand.
% Sources: N=natural-scope leverage diagnostic (validation), G=ME-GA record / pair counts,
% P=production docs, V=v3 scope declaration, S=v3 search declaration, B=benchmark summary (--bench), PENDING=not yet produced.
\providecommand{\PENDING}[1]{\textcolor{red}{\textbf{[#1]}}}
\makeatletter
\newcommand{\res}[1]{\ifcsname res@#1\endcsname\csname res@#1\endcsname\else\PENDING{undefined:#1}\fi}
\expandafter\def\csname res@natApW64\endcsname{0.20}% N
\expandafter\def\csname res@natApH64\endcsname{0.21}% N
\expandafter\def\csname res@natApW128\endcsname{0.42}% N
\expandafter\def\csname res@natApH128\endcsname{0.21}% N
\expandafter\def\csname res@natApW256\endcsname{0.42}% N
\expandafter\def\csname res@natApH256\endcsname{0.45}% N
\expandafter\def\csname res@natFrames64\endcsname{864}% N
\expandafter\def\csname res@natZeroShare64\endcsname{49.1}% N
\expandafter\def\csname res@natZeroBoth64\endcsname{370}% N
\expandafter\def\csname res@natZeroTraj64\endcsname{23}% N
\expandafter\def\csname res@natTraj64\endcsname{54}% N
\expandafter\def\csname res@natOffMinusNoRis64\endcsname{+0.019}% N
\expandafter\def\csname res@natGainAllMean64\endcsname{0.026}% N
\expandafter\def\csname res@natGainAllMedian64\endcsname{0.000}% N
\expandafter\def\csname res@natGainNzMean64\endcsname{0.051}% N
\expandafter\def\csname res@natGainNzMedian64\endcsname{0.015}% N
\expandafter\def\csname res@natGainNzPninety64\endcsname{0.139}% N
\expandafter\def\csname res@natGainNzMax64\endcsname{0.674}% N
\expandafter\def\csname res@natGainGtOne64\endcsname{6.7}% N
\expandafter\def\csname res@natGainGtFive64\endcsname{0.6}% N
\expandafter\def\csname res@natRateOffMean64\endcsname{3.234}% N
\expandafter\def\csname res@natRateNoRisMean64\endcsname{3.215}% N
\expandafter\def\csname res@natRateBestMean64\endcsname{3.260}% N
\expandafter\def\csname res@natRateOnMean64\endcsname{3.181}% N
\expandafter\def\csname res@natRateBestOwnMean64\endcsname{3.253}% N
\expandafter\def\csname res@natIN64\endcsname{$-$6.9}% N
\expandafter\def\csname res@natPrecErr64\endcsname{0}% N
\expandafter\def\csname res@natRisDirBlkLos64\endcsname{$-$55.6}% N
\expandafter\def\csname res@natBoundBlkLos64\endcsname{$-$41.6}% N
\expandafter\def\csname res@natUFBlkLos64\endcsname{340}% N
\expandafter\def\csname res@natRisDirRefLos64\endcsname{$-$55.7}% N
\expandafter\def\csname res@natBoundRefLos64\endcsname{$-$43.2}% N
\expandafter\def\csname res@natUFRefLos64\endcsname{350}% N
\expandafter\def\csname res@natRisDirLosLos64\endcsname{$-$59.2}% N
\expandafter\def\csname res@natBoundLosLos64\endcsname{$-$40.8}% N
\expandafter\def\csname res@natUFLosLos64\endcsname{969}% N
\expandafter\def\csname res@natRisDirBlkRef64\endcsname{+19.6}% N
\expandafter\def\csname res@natBoundBlkRef64\endcsname{$-$41.9}% N
\expandafter\def\csname res@natUFBlkRef64\endcsname{18}% N
\expandafter\def\csname res@natHrNLosLos64\endcsname{136}% N
\expandafter\def\csname res@natHrMedLosLos64\endcsname{0.027}% N
\expandafter\def\csname res@natHrGtLosLos64\endcsname{17}% N
\expandafter\def\csname res@natHrNBlkLos64\endcsname{115}% N
\expandafter\def\csname res@natHrMedBlkLos64\endcsname{0.023}% N
\expandafter\def\csname res@natHrGtBlkLos64\endcsname{24}% N
\expandafter\def\csname res@natZeroMunich64\endcsname{56}% N
\expandafter\def\csname res@natZeroRotterdam64\endcsname{28}% N
\expandafter\def\csname res@natZeroMaastricht64\endcsname{64}% N
\expandafter\def\csname res@natFrames128\endcsname{864}% N
\expandafter\def\csname res@natZeroShare128\endcsname{49.1}% N
\expandafter\def\csname res@natZeroBoth128\endcsname{376}% N
\expandafter\def\csname res@natZeroTraj128\endcsname{23}% N
\expandafter\def\csname res@natTraj128\endcsname{54}% N
\expandafter\def\csname res@natOffMinusNoRis128\endcsname{+0.025}% N
\expandafter\def\csname res@natGainAllMean128\endcsname{0.035}% N
\expandafter\def\csname res@natGainAllMedian128\endcsname{0.000}% N
\expandafter\def\csname res@natGainNzMean128\endcsname{0.069}% N
\expandafter\def\csname res@natGainNzMedian128\endcsname{0.022}% N
\expandafter\def\csname res@natGainNzPninety128\endcsname{0.170}% N
\expandafter\def\csname res@natGainNzMax128\endcsname{0.831}% N
\expandafter\def\csname res@natGainGtOne128\endcsname{10.1}% N
\expandafter\def\csname res@natGainGtFive128\endcsname{1.2}% N
\expandafter\def\csname res@natRateOffMean128\endcsname{3.315}% N
\expandafter\def\csname res@natRateNoRisMean128\endcsname{3.290}% N
\expandafter\def\csname res@natRateBestMean128\endcsname{3.350}% N
\expandafter\def\csname res@natRateOnMean128\endcsname{3.240}% N
\expandafter\def\csname res@natRateBestOwnMean128\endcsname{3.344}% N
\expandafter\def\csname res@natIN128\endcsname{$-$2.6}% N
\expandafter\def\csname res@natPrecErr128\endcsname{0}% N
\expandafter\def\csname res@natRisDirBlkLos128\endcsname{$-$52.8}% N
\expandafter\def\csname res@natBoundBlkLos128\endcsname{$-$36.0}% N
\expandafter\def\csname res@natUFBlkLos128\endcsname{333}% N
\expandafter\def\csname res@natRisDirRefLos128\endcsname{$-$49.6}% N
\expandafter\def\csname res@natBoundRefLos128\endcsname{$-$36.5}% N
\expandafter\def\csname res@natUFRefLos128\endcsname{358}% N
\expandafter\def\csname res@natRisDirLosLos128\endcsname{$-$57.6}% N
\expandafter\def\csname res@natBoundLosLos128\endcsname{$-$34.2}% N
\expandafter\def\csname res@natUFLosLos128\endcsname{972}% N
\expandafter\def\csname res@natRisDirBlkRef128\endcsname{+19.4}% N
\expandafter\def\csname res@natBoundBlkRef128\endcsname{$-$35.2}% N
\expandafter\def\csname res@natUFBlkRef128\endcsname{18}% N
\expandafter\def\csname res@natHrNLosLos128\endcsname{118}% N
\expandafter\def\csname res@natHrMedLosLos128\endcsname{0.050}% N
\expandafter\def\csname res@natHrGtLosLos128\endcsname{33}% N
\expandafter\def\csname res@natHrNBlkLos128\endcsname{104}% N
\expandafter\def\csname res@natHrMedBlkLos128\endcsname{0.029}% N
\expandafter\def\csname res@natHrGtBlkLos128\endcsname{27}% N
\expandafter\def\csname res@natZeroMunich128\endcsname{56}% N
\expandafter\def\csname res@natZeroRotterdam128\endcsname{28}% N
\expandafter\def\csname res@natZeroMaastricht128\endcsname{64}% N
\expandafter\def\csname res@natFrames256\endcsname{864}% N
\expandafter\def\csname res@natZeroShare256\endcsname{49.7}% N
\expandafter\def\csname res@natZeroBoth256\endcsname{381}% N
\expandafter\def\csname res@natZeroTraj256\endcsname{23}% N
\expandafter\def\csname res@natTraj256\endcsname{54}% N
\expandafter\def\csname res@natOffMinusNoRis256\endcsname{+0.043}% N
\expandafter\def\csname res@natGainAllMean256\endcsname{0.060}% N
\expandafter\def\csname res@natGainAllMedian256\endcsname{0.000}% N
\expandafter\def\csname res@natGainNzMean256\endcsname{0.120}% N
\expandafter\def\csname res@natGainNzMedian256\endcsname{0.045}% N
\expandafter\def\csname res@natGainNzPninety256\endcsname{0.292}% N
\expandafter\def\csname res@natGainNzMax256\endcsname{1.488}% N
\expandafter\def\csname res@natGainGtOne256\endcsname{16.3}% N
\expandafter\def\csname res@natGainGtFive256\endcsname{2.8}% N
\expandafter\def\csname res@natRateOffMean256\endcsname{3.242}% N
\expandafter\def\csname res@natRateNoRisMean256\endcsname{3.199}% N
\expandafter\def\csname res@natRateBestMean256\endcsname{3.303}% N
\expandafter\def\csname res@natRateOnMean256\endcsname{3.137}% N
\expandafter\def\csname res@natRateBestOwnMean256\endcsname{3.293}% N
\expandafter\def\csname res@natIN256\endcsname{$-$7.2}% N
\expandafter\def\csname res@natPrecErr256\endcsname{0}% N
\expandafter\def\csname res@natRisDirBlkLos256\endcsname{$-$40.9}% N
\expandafter\def\csname res@natBoundBlkLos256\endcsname{$-$29.4}% N
\expandafter\def\csname res@natUFBlkLos256\endcsname{345}% N
\expandafter\def\csname res@natRisDirRefLos256\endcsname{$-$43.8}% N
\expandafter\def\csname res@natBoundRefLos256\endcsname{$-$29.2}% N
\expandafter\def\csname res@natUFRefLos256\endcsname{346}% N
\expandafter\def\csname res@natRisDirLosLos256\endcsname{$-$56.0}% N
\expandafter\def\csname res@natBoundLosLos256\endcsname{$-$27.6}% N
\expandafter\def\csname res@natUFLosLos256\endcsname{972}% N
\expandafter\def\csname res@natRisDirBlkRef256\endcsname{+19.6}% N
\expandafter\def\csname res@natBoundBlkRef256\endcsname{$-$28.6}% N
\expandafter\def\csname res@natUFBlkRef256\endcsname{18}% N
\expandafter\def\csname res@natHrNLosLos256\endcsname{110}% N
\expandafter\def\csname res@natHrMedLosLos256\endcsname{0.125}% N
\expandafter\def\csname res@natHrGtLosLos256\endcsname{55}% N
\expandafter\def\csname res@natHrNBlkLos256\endcsname{140}% N
\expandafter\def\csname res@natHrMedBlkLos256\endcsname{0.047}% N
\expandafter\def\csname res@natHrGtBlkLos256\endcsname{35}% N
\expandafter\def\csname res@natZeroMunich256\endcsname{56}% N
\expandafter\def\csname res@natZeroRotterdam256\endcsname{28}% N
\expandafter\def\csname res@natZeroMaastricht256\endcsname{65}% N
\expandafter\def\csname res@natFramesTotal\endcsname{2,592}% N
\expandafter\def\csname res@freqGHz\endcsname{5.53}% N
\expandafter\def\csname res@lambdaMm\endcsname{54.2}% N
\expandafter\def\csname res@boundXtenDb\endcsname{11}% N
\expandafter\def\csname res@natDbsrisMedian\endcsname{112}% N
\expandafter\def\csname res@natDbsrisMax\endcsname{150}% N
\expandafter\def\csname res@natDrisUserLo\endcsname{63}% N
\expandafter\def\csname res@natDrisUserHi\endcsname{136}% N
\expandafter\def\csname res@prodTrajPerMap\endcsname{120}% P
\expandafter\def\csname res@prodTraj\endcsname{360}% P
\expandafter\def\csname res@prodFramesPerTraj\endcsname{16}% P
\expandafter\def\csname res@prodFrames\endcsname{5,760}% P
\expandafter\def\csname res@prodArrayFrames\endcsname{17,280}% P
\expandafter\def\csname res@prodValTraj\endcsname{54}% P
\expandafter\def\csname res@siteDminNat\endcsname{20}% P
\expandafter\def\csname res@siteDmaxNat\endcsname{150}% P
\expandafter\def\csname res@negSlotProb\endcsname{0.15}% P
\expandafter\def\csname res@rtSamples\endcsname{10,000}% P
\expandafter\def\csname res@rtDepth\endcsname{2}% P
\expandafter\def\csname res@openGapPct\endcsname{4.7}% P
\expandafter\def\csname res@megaPop\endcsname{100}% G
\expandafter\def\csname res@megaIter\endcsname{100}% G
\expandafter\def\csname res@megaDecisions\endcsname{24}% G
\expandafter\def\csname res@megaSensItems\endcsname{9}% G
\expandafter\def\csname res@nPairsSixtyFour\endcsname{2,016}% G
\expandafter\def\csname res@nPairsOneTwentyEight\endcsname{8,128}% G
\expandafter\def\csname res@nPairsTwoFiftySix\endcsname{32,640}% G
\expandafter\def\csname res@megaWoodM\endcsname{24}% G
\expandafter\def\csname res@megaWoodMin\endcsname{12}% G
\expandafter\def\csname res@nsgaPop\endcsname{20}% G
\expandafter\def\csname res@nsgaPc\endcsname{0.9}% G
\expandafter\def\csname res@nsgaGridP\endcsname{20/50/100}% G
\expandafter\def\csname res@nsgaGridPc\endcsname{0.7/0.9}% G
\expandafter\def\csname res@nsgaGridN\endcsname{6}% G
\expandafter\def\csname res@vMap\endcsname{Rotterdam}% V
\expandafter\def\csname res@vTraj\endcsname{120}% V
\expandafter\def\csname res@vFrames\endcsname{16}% V
\expandafter\def\csname res@vTrainTraj\endcsname{84}% V
\expandafter\def\csname res@vValTraj\endcsname{18}% V
\expandafter\def\csname res@vTestTraj\endcsname{18}% V
\expandafter\def\csname res@vDmin\endcsname{5}% V
\expandafter\def\csname res@vDmax\endcsname{20}% V
\expandafter\def\csname res@vDOneMin\endcsname{7.8}% V
\expandafter\def\csname res@vDOneMed\endcsname{13.7}% V
\expandafter\def\csname res@vDOneMax\endcsname{18.7}% V
\expandafter\def\csname res@vSites\endcsname{4}% V
\expandafter\def\csname res@vDTwoLo\endcsname{26.7}% V
\expandafter\def\csname res@vDTwoMed\endcsname{32.3}% V
\expandafter\def\csname res@vDTwoHi\endcsname{36.4}% V
\expandafter\def\csname res@vRadius\endcsname{40}% V
\expandafter\def\csname res@vAngle\endcsname{80}% V
\expandafter\def\csname res@vServedUF\endcsname{4,800}% V
\expandafter\def\csname res@vSamples\endcsname{10^{6}}% V
\expandafter\def\csname res@vPathCap\endcsname{2\times10^{6}}% V
\expandafter\def\csname res@vDepth\endcsname{2}% V
\expandafter\def\csname res@vUnitsOk\endcsname{357}% V
\expandafter\def\csname res@vUnits\endcsname{360}% V
\expandafter\def\csname res@vIncluded\endcsname{119}% V
\expandafter\def\csname res@vGainMed64\endcsname{0.302}% V
\expandafter\def\csname res@vGainMed128\endcsname{0.418}% V
\expandafter\def\csname res@vGainMed256\endcsname{0.782}% V
\expandafter\def\csname res@vGainMean64\endcsname{0.450}% V
\expandafter\def\csname res@vGainPninety64\endcsname{1.089}% V
\expandafter\def\csname res@vGainMean128\endcsname{0.534}% V
\expandafter\def\csname res@vGainPninety128\endcsname{1.379}% V
\expandafter\def\csname res@vGainMean256\endcsname{1.097}% V
\expandafter\def\csname res@vGainPninety256\endcsname{2.760}% V
\expandafter\def\csname res@vPathless\endcsname{0.0}% V
\expandafter\def\csname res@vSeeds\endcsname{10}% S
\expandafter\def\csname res@vBoot\endcsname{10,000}% S
\expandafter\def\csname res@vFamily\endcsname{five}% S
\expandafter\def\csname res@vMaxMiss\endcsname{0.5}% S
\expandafter\def\csname res@vRateGreedyFive64\endcsname{\PENDING{vRateGreedyFive64}}% PENDING
\expandafter\def\csname res@vRateGreedyTen64\endcsname{\PENDING{vRateGreedyTen64}}% PENDING
\expandafter\def\csname res@vRateGreedyFifty64\endcsname{\PENDING{vRateGreedyFifty64}}% PENDING
\expandafter\def\csname res@vRateGreedyHundred64\endcsname{\PENDING{vRateGreedyHundred64}}% PENDING
\expandafter\def\csname res@vLatGreedy64\endcsname{\PENDING{vLatGreedy64}}% PENDING
\expandafter\def\csname res@vLatPGreedy64\endcsname{\PENDING{vLatPGreedy64}}% PENDING
\expandafter\def\csname res@vEvalsGreedy64\endcsname{\PENDING{vEvalsGreedy64}}% PENDING
\expandafter\def\csname res@vMissGreedy64\endcsname{\PENDING{vMissGreedy64}}% PENDING
\expandafter\def\csname res@vRatePairFive64\endcsname{\PENDING{vRatePairFive64}}% PENDING
\expandafter\def\csname res@vRatePairTen64\endcsname{\PENDING{vRatePairTen64}}% PENDING
\expandafter\def\csname res@vRatePairFifty64\endcsname{\PENDING{vRatePairFifty64}}% PENDING
\expandafter\def\csname res@vRatePairHundred64\endcsname{\PENDING{vRatePairHundred64}}% PENDING
\expandafter\def\csname res@vLatPair64\endcsname{\PENDING{vLatPair64}}% PENDING
\expandafter\def\csname res@vLatPPair64\endcsname{\PENDING{vLatPPair64}}% PENDING
\expandafter\def\csname res@vEvalsPair64\endcsname{\PENDING{vEvalsPair64}}% PENDING
\expandafter\def\csname res@vMissPair64\endcsname{\PENDING{vMissPair64}}% PENDING
\expandafter\def\csname res@vRateBthreeFive64\endcsname{\PENDING{vRateBthreeFive64}}% PENDING
\expandafter\def\csname res@vRateBthreeTen64\endcsname{\PENDING{vRateBthreeTen64}}% PENDING
\expandafter\def\csname res@vRateBthreeFifty64\endcsname{\PENDING{vRateBthreeFifty64}}% PENDING
\expandafter\def\csname res@vRateBthreeHundred64\endcsname{\PENDING{vRateBthreeHundred64}}% PENDING
\expandafter\def\csname res@vLatBthree64\endcsname{\PENDING{vLatBthree64}}% PENDING
\expandafter\def\csname res@vLatPBthree64\endcsname{\PENDING{vLatPBthree64}}% PENDING
\expandafter\def\csname res@vEvalsBthree64\endcsname{\PENDING{vEvalsBthree64}}% PENDING
\expandafter\def\csname res@vMissBthree64\endcsname{\PENDING{vMissBthree64}}% PENDING
\expandafter\def\csname res@vRateBfiveFive64\endcsname{\PENDING{vRateBfiveFive64}}% PENDING
\expandafter\def\csname res@vRateBfiveTen64\endcsname{\PENDING{vRateBfiveTen64}}% PENDING
\expandafter\def\csname res@vRateBfiveFifty64\endcsname{\PENDING{vRateBfiveFifty64}}% PENDING
\expandafter\def\csname res@vRateBfiveHundred64\endcsname{\PENDING{vRateBfiveHundred64}}% PENDING
\expandafter\def\csname res@vLatBfive64\endcsname{\PENDING{vLatBfive64}}% PENDING
\expandafter\def\csname res@vLatPBfive64\endcsname{\PENDING{vLatPBfive64}}% PENDING
\expandafter\def\csname res@vEvalsBfive64\endcsname{\PENDING{vEvalsBfive64}}% PENDING
\expandafter\def\csname res@vMissBfive64\endcsname{\PENDING{vMissBfive64}}% PENDING
\expandafter\def\csname res@vRateBtsFive64\endcsname{\PENDING{vRateBtsFive64}}% PENDING
\expandafter\def\csname res@vRateBtsTen64\endcsname{\PENDING{vRateBtsTen64}}% PENDING
\expandafter\def\csname res@vRateBtsFifty64\endcsname{\PENDING{vRateBtsFifty64}}% PENDING
\expandafter\def\csname res@vRateBtsHundred64\endcsname{\PENDING{vRateBtsHundred64}}% PENDING
\expandafter\def\csname res@vLatBts64\endcsname{\PENDING{vLatBts64}}% PENDING
\expandafter\def\csname res@vLatPBts64\endcsname{\PENDING{vLatPBts64}}% PENDING
\expandafter\def\csname res@vEvalsBts64\endcsname{\PENDING{vEvalsBts64}}% PENDING
\expandafter\def\csname res@vMissBts64\endcsname{\PENDING{vMissBts64}}% PENDING
\expandafter\def\csname res@vRateMegaFive64\endcsname{\PENDING{vRateMegaFive64}}% PENDING
\expandafter\def\csname res@vRateMegaTen64\endcsname{\PENDING{vRateMegaTen64}}% PENDING
\expandafter\def\csname res@vRateMegaFifty64\endcsname{\PENDING{vRateMegaFifty64}}% PENDING
\expandafter\def\csname res@vRateMegaHundred64\endcsname{\PENDING{vRateMegaHundred64}}% PENDING
\expandafter\def\csname res@vLatMega64\endcsname{\PENDING{vLatMega64}}% PENDING
\expandafter\def\csname res@vLatPMega64\endcsname{\PENDING{vLatPMega64}}% PENDING
\expandafter\def\csname res@vEvalsMega64\endcsname{\PENDING{vEvalsMega64}}% PENDING
\expandafter\def\csname res@vMissMega64\endcsname{\PENDING{vMissMega64}}% PENDING
\expandafter\def\csname res@vRateNsgaFive64\endcsname{\PENDING{vRateNsgaFive64}}% PENDING
\expandafter\def\csname res@vRateNsgaTen64\endcsname{\PENDING{vRateNsgaTen64}}% PENDING
\expandafter\def\csname res@vRateNsgaFifty64\endcsname{\PENDING{vRateNsgaFifty64}}% PENDING
\expandafter\def\csname res@vRateNsgaHundred64\endcsname{\PENDING{vRateNsgaHundred64}}% PENDING
\expandafter\def\csname res@vLatNsga64\endcsname{\PENDING{vLatNsga64}}% PENDING
\expandafter\def\csname res@vLatPNsga64\endcsname{\PENDING{vLatPNsga64}}% PENDING
\expandafter\def\csname res@vEvalsNsga64\endcsname{\PENDING{vEvalsNsga64}}% PENDING
\expandafter\def\csname res@vMissNsga64\endcsname{\PENDING{vMissNsga64}}% PENDING
\expandafter\def\csname res@vHOne64\endcsname{\PENDING{vHOne64}}% PENDING
\expandafter\def\csname res@vHOneVerdict64\endcsname{\PENDING{vHOneVerdict64}}% PENDING
\expandafter\def\csname res@vHTwo64\endcsname{\PENDING{vHTwo64}}% PENDING
\expandafter\def\csname res@vHTwoVerdict64\endcsname{\PENDING{vHTwoVerdict64}}% PENDING
\expandafter\def\csname res@vHFive64\endcsname{\PENDING{vHFive64}}% PENDING
\expandafter\def\csname res@vHFiveVerdict64\endcsname{\PENDING{vHFiveVerdict64}}% PENDING
\expandafter\def\csname res@vRateGreedyFive128\endcsname{\PENDING{vRateGreedyFive128}}% PENDING
\expandafter\def\csname res@vRateGreedyTen128\endcsname{\PENDING{vRateGreedyTen128}}% PENDING
\expandafter\def\csname res@vRateGreedyFifty128\endcsname{\PENDING{vRateGreedyFifty128}}% PENDING
\expandafter\def\csname res@vRateGreedyHundred128\endcsname{\PENDING{vRateGreedyHundred128}}% PENDING
\expandafter\def\csname res@vLatGreedy128\endcsname{\PENDING{vLatGreedy128}}% PENDING
\expandafter\def\csname res@vLatPGreedy128\endcsname{\PENDING{vLatPGreedy128}}% PENDING
\expandafter\def\csname res@vEvalsGreedy128\endcsname{\PENDING{vEvalsGreedy128}}% PENDING
\expandafter\def\csname res@vMissGreedy128\endcsname{\PENDING{vMissGreedy128}}% PENDING
\expandafter\def\csname res@vRatePairFive128\endcsname{\PENDING{vRatePairFive128}}% PENDING
\expandafter\def\csname res@vRatePairTen128\endcsname{\PENDING{vRatePairTen128}}% PENDING
\expandafter\def\csname res@vRatePairFifty128\endcsname{\PENDING{vRatePairFifty128}}% PENDING
\expandafter\def\csname res@vRatePairHundred128\endcsname{\PENDING{vRatePairHundred128}}% PENDING
\expandafter\def\csname res@vLatPair128\endcsname{\PENDING{vLatPair128}}% PENDING
\expandafter\def\csname res@vLatPPair128\endcsname{\PENDING{vLatPPair128}}% PENDING
\expandafter\def\csname res@vEvalsPair128\endcsname{\PENDING{vEvalsPair128}}% PENDING
\expandafter\def\csname res@vMissPair128\endcsname{\PENDING{vMissPair128}}% PENDING
\expandafter\def\csname res@vRateBthreeFive128\endcsname{\PENDING{vRateBthreeFive128}}% PENDING
\expandafter\def\csname res@vRateBthreeTen128\endcsname{\PENDING{vRateBthreeTen128}}% PENDING
\expandafter\def\csname res@vRateBthreeFifty128\endcsname{\PENDING{vRateBthreeFifty128}}% PENDING
\expandafter\def\csname res@vRateBthreeHundred128\endcsname{\PENDING{vRateBthreeHundred128}}% PENDING
\expandafter\def\csname res@vLatBthree128\endcsname{\PENDING{vLatBthree128}}% PENDING
\expandafter\def\csname res@vLatPBthree128\endcsname{\PENDING{vLatPBthree128}}% PENDING
\expandafter\def\csname res@vEvalsBthree128\endcsname{\PENDING{vEvalsBthree128}}% PENDING
\expandafter\def\csname res@vMissBthree128\endcsname{\PENDING{vMissBthree128}}% PENDING
\expandafter\def\csname res@vRateBfiveFive128\endcsname{\PENDING{vRateBfiveFive128}}% PENDING
\expandafter\def\csname res@vRateBfiveTen128\endcsname{\PENDING{vRateBfiveTen128}}% PENDING
\expandafter\def\csname res@vRateBfiveFifty128\endcsname{\PENDING{vRateBfiveFifty128}}% PENDING
\expandafter\def\csname res@vRateBfiveHundred128\endcsname{\PENDING{vRateBfiveHundred128}}% PENDING
\expandafter\def\csname res@vLatBfive128\endcsname{\PENDING{vLatBfive128}}% PENDING
\expandafter\def\csname res@vLatPBfive128\endcsname{\PENDING{vLatPBfive128}}% PENDING
\expandafter\def\csname res@vEvalsBfive128\endcsname{\PENDING{vEvalsBfive128}}% PENDING
\expandafter\def\csname res@vMissBfive128\endcsname{\PENDING{vMissBfive128}}% PENDING
\expandafter\def\csname res@vRateBtsFive128\endcsname{\PENDING{vRateBtsFive128}}% PENDING
\expandafter\def\csname res@vRateBtsTen128\endcsname{\PENDING{vRateBtsTen128}}% PENDING
\expandafter\def\csname res@vRateBtsFifty128\endcsname{\PENDING{vRateBtsFifty128}}% PENDING
\expandafter\def\csname res@vRateBtsHundred128\endcsname{\PENDING{vRateBtsHundred128}}% PENDING
\expandafter\def\csname res@vLatBts128\endcsname{\PENDING{vLatBts128}}% PENDING
\expandafter\def\csname res@vLatPBts128\endcsname{\PENDING{vLatPBts128}}% PENDING
\expandafter\def\csname res@vEvalsBts128\endcsname{\PENDING{vEvalsBts128}}% PENDING
\expandafter\def\csname res@vMissBts128\endcsname{\PENDING{vMissBts128}}% PENDING
\expandafter\def\csname res@vRateMegaFive128\endcsname{\PENDING{vRateMegaFive128}}% PENDING
\expandafter\def\csname res@vRateMegaTen128\endcsname{\PENDING{vRateMegaTen128}}% PENDING
\expandafter\def\csname res@vRateMegaFifty128\endcsname{\PENDING{vRateMegaFifty128}}% PENDING
\expandafter\def\csname res@vRateMegaHundred128\endcsname{\PENDING{vRateMegaHundred128}}% PENDING
\expandafter\def\csname res@vLatMega128\endcsname{\PENDING{vLatMega128}}% PENDING
\expandafter\def\csname res@vLatPMega128\endcsname{\PENDING{vLatPMega128}}% PENDING
\expandafter\def\csname res@vEvalsMega128\endcsname{\PENDING{vEvalsMega128}}% PENDING
\expandafter\def\csname res@vMissMega128\endcsname{\PENDING{vMissMega128}}% PENDING
\expandafter\def\csname res@vRateNsgaFive128\endcsname{\PENDING{vRateNsgaFive128}}% PENDING
\expandafter\def\csname res@vRateNsgaTen128\endcsname{\PENDING{vRateNsgaTen128}}% PENDING
\expandafter\def\csname res@vRateNsgaFifty128\endcsname{\PENDING{vRateNsgaFifty128}}% PENDING
\expandafter\def\csname res@vRateNsgaHundred128\endcsname{\PENDING{vRateNsgaHundred128}}% PENDING
\expandafter\def\csname res@vLatNsga128\endcsname{\PENDING{vLatNsga128}}% PENDING
\expandafter\def\csname res@vLatPNsga128\endcsname{\PENDING{vLatPNsga128}}% PENDING
\expandafter\def\csname res@vEvalsNsga128\endcsname{\PENDING{vEvalsNsga128}}% PENDING
\expandafter\def\csname res@vMissNsga128\endcsname{\PENDING{vMissNsga128}}% PENDING
\expandafter\def\csname res@vHOne128\endcsname{\PENDING{vHOne128}}% PENDING
\expandafter\def\csname res@vHOneVerdict128\endcsname{\PENDING{vHOneVerdict128}}% PENDING
\expandafter\def\csname res@vHTwo128\endcsname{\PENDING{vHTwo128}}% PENDING
\expandafter\def\csname res@vHTwoVerdict128\endcsname{\PENDING{vHTwoVerdict128}}% PENDING
\expandafter\def\csname res@vHFive128\endcsname{\PENDING{vHFive128}}% PENDING
\expandafter\def\csname res@vHFiveVerdict128\endcsname{\PENDING{vHFiveVerdict128}}% PENDING
\expandafter\def\csname res@vRateGreedyFive256\endcsname{\PENDING{vRateGreedyFive256}}% PENDING
\expandafter\def\csname res@vRateGreedyTen256\endcsname{\PENDING{vRateGreedyTen256}}% PENDING
\expandafter\def\csname res@vRateGreedyFifty256\endcsname{\PENDING{vRateGreedyFifty256}}% PENDING
\expandafter\def\csname res@vRateGreedyHundred256\endcsname{\PENDING{vRateGreedyHundred256}}% PENDING
\expandafter\def\csname res@vLatGreedy256\endcsname{\PENDING{vLatGreedy256}}% PENDING
\expandafter\def\csname res@vLatPGreedy256\endcsname{\PENDING{vLatPGreedy256}}% PENDING
\expandafter\def\csname res@vEvalsGreedy256\endcsname{\PENDING{vEvalsGreedy256}}% PENDING
\expandafter\def\csname res@vMissGreedy256\endcsname{\PENDING{vMissGreedy256}}% PENDING
\expandafter\def\csname res@vRatePairFive256\endcsname{\PENDING{vRatePairFive256}}% PENDING
\expandafter\def\csname res@vRatePairTen256\endcsname{\PENDING{vRatePairTen256}}% PENDING
\expandafter\def\csname res@vRatePairFifty256\endcsname{\PENDING{vRatePairFifty256}}% PENDING
\expandafter\def\csname res@vRatePairHundred256\endcsname{\PENDING{vRatePairHundred256}}% PENDING
\expandafter\def\csname res@vLatPair256\endcsname{\PENDING{vLatPair256}}% PENDING
\expandafter\def\csname res@vLatPPair256\endcsname{\PENDING{vLatPPair256}}% PENDING
\expandafter\def\csname res@vEvalsPair256\endcsname{\PENDING{vEvalsPair256}}% PENDING
\expandafter\def\csname res@vMissPair256\endcsname{\PENDING{vMissPair256}}% PENDING
\expandafter\def\csname res@vRateBthreeFive256\endcsname{\PENDING{vRateBthreeFive256}}% PENDING
\expandafter\def\csname res@vRateBthreeTen256\endcsname{\PENDING{vRateBthreeTen256}}% PENDING
\expandafter\def\csname res@vRateBthreeFifty256\endcsname{\PENDING{vRateBthreeFifty256}}% PENDING
\expandafter\def\csname res@vRateBthreeHundred256\endcsname{\PENDING{vRateBthreeHundred256}}% PENDING
\expandafter\def\csname res@vLatBthree256\endcsname{\PENDING{vLatBthree256}}% PENDING
\expandafter\def\csname res@vLatPBthree256\endcsname{\PENDING{vLatPBthree256}}% PENDING
\expandafter\def\csname res@vEvalsBthree256\endcsname{\PENDING{vEvalsBthree256}}% PENDING
\expandafter\def\csname res@vMissBthree256\endcsname{\PENDING{vMissBthree256}}% PENDING
\expandafter\def\csname res@vRateBfiveFive256\endcsname{\PENDING{vRateBfiveFive256}}% PENDING
\expandafter\def\csname res@vRateBfiveTen256\endcsname{\PENDING{vRateBfiveTen256}}% PENDING
\expandafter\def\csname res@vRateBfiveFifty256\endcsname{\PENDING{vRateBfiveFifty256}}% PENDING
\expandafter\def\csname res@vRateBfiveHundred256\endcsname{\PENDING{vRateBfiveHundred256}}% PENDING
\expandafter\def\csname res@vLatBfive256\endcsname{\PENDING{vLatBfive256}}% PENDING
\expandafter\def\csname res@vLatPBfive256\endcsname{\PENDING{vLatPBfive256}}% PENDING
\expandafter\def\csname res@vEvalsBfive256\endcsname{\PENDING{vEvalsBfive256}}% PENDING
\expandafter\def\csname res@vMissBfive256\endcsname{\PENDING{vMissBfive256}}% PENDING
\expandafter\def\csname res@vRateBtsFive256\endcsname{\PENDING{vRateBtsFive256}}% PENDING
\expandafter\def\csname res@vRateBtsTen256\endcsname{\PENDING{vRateBtsTen256}}% PENDING
\expandafter\def\csname res@vRateBtsFifty256\endcsname{\PENDING{vRateBtsFifty256}}% PENDING
\expandafter\def\csname res@vRateBtsHundred256\endcsname{\PENDING{vRateBtsHundred256}}% PENDING
\expandafter\def\csname res@vLatBts256\endcsname{\PENDING{vLatBts256}}% PENDING
\expandafter\def\csname res@vLatPBts256\endcsname{\PENDING{vLatPBts256}}% PENDING
\expandafter\def\csname res@vEvalsBts256\endcsname{\PENDING{vEvalsBts256}}% PENDING
\expandafter\def\csname res@vMissBts256\endcsname{\PENDING{vMissBts256}}% PENDING
\expandafter\def\csname res@vRateMegaFive256\endcsname{\PENDING{vRateMegaFive256}}% PENDING
\expandafter\def\csname res@vRateMegaTen256\endcsname{\PENDING{vRateMegaTen256}}% PENDING
\expandafter\def\csname res@vRateMegaFifty256\endcsname{\PENDING{vRateMegaFifty256}}% PENDING
\expandafter\def\csname res@vRateMegaHundred256\endcsname{\PENDING{vRateMegaHundred256}}% PENDING
\expandafter\def\csname res@vLatMega256\endcsname{\PENDING{vLatMega256}}% PENDING
\expandafter\def\csname res@vLatPMega256\endcsname{\PENDING{vLatPMega256}}% PENDING
\expandafter\def\csname res@vEvalsMega256\endcsname{\PENDING{vEvalsMega256}}% PENDING
\expandafter\def\csname res@vMissMega256\endcsname{\PENDING{vMissMega256}}% PENDING
\expandafter\def\csname res@vRateNsgaFive256\endcsname{\PENDING{vRateNsgaFive256}}% PENDING
\expandafter\def\csname res@vRateNsgaTen256\endcsname{\PENDING{vRateNsgaTen256}}% PENDING
\expandafter\def\csname res@vRateNsgaFifty256\endcsname{\PENDING{vRateNsgaFifty256}}% PENDING
\expandafter\def\csname res@vRateNsgaHundred256\endcsname{\PENDING{vRateNsgaHundred256}}% PENDING
\expandafter\def\csname res@vLatNsga256\endcsname{\PENDING{vLatNsga256}}% PENDING
\expandafter\def\csname res@vLatPNsga256\endcsname{\PENDING{vLatPNsga256}}% PENDING
\expandafter\def\csname res@vEvalsNsga256\endcsname{\PENDING{vEvalsNsga256}}% PENDING
\expandafter\def\csname res@vMissNsga256\endcsname{\PENDING{vMissNsga256}}% PENDING
\expandafter\def\csname res@vHOne256\endcsname{\PENDING{vHOne256}}% PENDING
\expandafter\def\csname res@vHOneVerdict256\endcsname{\PENDING{vHOneVerdict256}}% PENDING
\expandafter\def\csname res@vHTwo256\endcsname{\PENDING{vHTwo256}}% PENDING
\expandafter\def\csname res@vHTwoVerdict256\endcsname{\PENDING{vHTwoVerdict256}}% PENDING
\expandafter\def\csname res@vHFive256\endcsname{\PENDING{vHFive256}}% PENDING
\expandafter\def\csname res@vHFiveVerdict256\endcsname{\PENDING{vHFiveVerdict256}}% PENDING
\expandafter\def\csname res@vBenchLoaded\endcsname{no}% B
\makeatother

\newcommand{\bts}{\textsc{b35}}
\newcommand{\Hd}{H_{\mathrm{d}}}
\renewcommand{\algorithmiccomment}[1]{\hfill{\footnotesize$\triangleright$ #1}}

\begin{document}
\title{Title TBD}
\titlerunning{Title TBD}
\author{Author One\inst{1}\orcidID{0000-0000-0000-0000} \and
Author Two\inst{1} \and
Author Three\inst{2}}
\authorrunning{A. One et al.}
\institute{Department, University, City, Country\\
\email{\{author.one,author.two\}@example.edu} \and
Department, Institute, City, Country\\
\email{author.three@example.org}}
\maketitle

\begin{abstract}
Configuring a 1-bit reconfigurable intelligent surface (RIS) to maximise the worst user's rate is a binary max--min problem that must be solved again whenever users move, within tens of milliseconds. Evolutionary controllers are a common answer, but on a mutually coupled surface each candidate needs a coupled-network evaluation. We study \bts{}, a memetic iterated local search that scores every single and every pair flip exactly through rank-one and rank-two Woodbury updates of the coupled multiport network, restarts from two predeclared ON densities, and every third restart recombines two cached local optima, all inside a hard CPU deadline that includes set-up. Its ingredients are known individually; the contribution is their combination for a multiuser max--min rate objective on a coupled surface under a complete-runtime deadline, and a deadline-fair evaluation. We evaluate it on GeoRIS-Bench, a city-scale benchmark that couples ray tracing over open LoD2 building models with a full method-of-moments model of 64-, 128- and 256-element surfaces, against single-flip and exact pair search, two ablations, our reimplementation of the memory-enhanced genetic algorithm ME-GA, and NSGA-II, on worst-user spectral efficiency, optimisation latency and computational budget at deadlines of 5--100\,ms.
\keywords{Reconfigurable intelligent surface \and Combinatorial optimisation \and Anytime search \and Local search \and Evolutionary algorithms \and Mutual coupling.}
\end{abstract}

%======================================================================
\section{Introduction}
%======================================================================
A reconfigurable intelligent surface is a passive array whose elements switch between a few load states, so that the field it re-radiates can be shaped without active transmit chains~\cite{li2026tutorial}. With 1-bit elements, configuring an $N$-element surface means choosing one of $2^N$ binary states. When the surface serves several users from a multi-antenna base station (BS), a natural objective is the rate of the worst served user. This objective is non-smooth, and once mutual coupling between elements is modelled~\cite{qian2021mutual,zheng2024mutual,nerini2024universal,nerini2026global}, it is also non-separable: the effect of one element depends on the state of all others. Users move, so the problem must be re-solved every few tens of milliseconds, on the hardware that controls the surface~\cite{urbinati2026ondevice}. The practical question is therefore not which method finds the global optimum, but which method returns the best state \emph{by a fixed deadline}.

Evolutionary algorithms are a common choice for this setting~\cite{stylianopoulos2022dlga,zardi2024memory}. Zardi, Oliveri and Massa~\cite{zardi2024memory} proposed ME-GA, a genetic algorithm whose memory of earlier good individuals lets it track a changing channel, and NSGA-II~\cite{deb2002fast} is the standard elitist evolutionary algorithm. Under a deadline their difficulty is cost: each individual needs its own coupled-network evaluation, a linear solve with an $N\times N$ impedance matrix unless it is close to an already evaluated state. Single-flip local search is cheaper, because a flip changes one diagonal load and can be scored by a rank-one update, but it stalls: a state can be optimal with respect to every single flip and still be improved by flipping two coupled elements together.

This paper builds a deadline-bounded search on top of exact pair moves. Pair flips, restarts, elite recombination and low-rank updates are each established~\cite{merz2002bqp,lourenco2010ils,merz2004memetic,prodhomme2023efficient}, and to our knowledge prior RIS benchmarks stop at single-flip descent~\cite{hammami2026statistical}. Our contributions are:
\begin{enumerate}
\item \textbf{\bts{}}, a memetic iterated local search for coupled 1-bit RIS control that combines these ingredients for a non-quadratic multiuser max--min rate objective. It scores all $N$ single flips and all $N(N-1)/2$ pair flips exactly with rank-one and rank-two updates of the coupled network, and diversifies through two predeclared restart densities and elite recombination of cached local optima (Sect.~\ref{sec:method}).
\item \textbf{A deadline-fair evaluation protocol} with three metrics shared by all methods -- worst-user spectral efficiency, optimisation latency and computational budget -- in which the clock covers evaluator set-up and every scan, the decision rules are frozen before a single held-out run, and every method faces the same evaluator and deadline (Sect.~\ref{sec:protocol}).
\item \textbf{A city-scale evaluation} on GeoRIS-Bench~\cite{georisbench}, built from LoD2 building models, ray tracing and a full method-of-moments surface model, against single-flip and exact pair search, the two ablations, ME-GA and NSGA-II at 64, 128 and 256 elements (Sect.~\ref{sec:results}).
\end{enumerate}

%======================================================================
\section{Related Work}\label{sec:related}
%======================================================================
\textbf{Learning-based configuration.} Deep networks that map channel observations to RIS states have been implemented on FPGAs~\cite{martin2024hardware,padial2026accelerating} and on embedded devices, where the latency--rate trade-off is explicit~\cite{urbinati2026ondevice}. Other works add robustness to practical constraints~\cite{xu2022robust}, account for coupling and partial channel knowledge~\cite{peng2025risnet}, or train generative surrogates on full-wave data~\cite{liu2026physics}. These methods amortise optimisation into training; the present work instead studies search at decision time.

\textbf{Evolutionary and discrete search.} Genetic algorithms are widely used for discrete RIS configuration~\cite{stylianopoulos2022dlga,radar2024onebitga}. ME-GA~\cite{zardi2024memory} adapts its mutation and crossover rates to population diversity and uses a memory of past best individuals to react to channel changes; it was not benchmarked against local search or under a deadline. NSGA-II~\cite{deb2002fast} has been used for binary metasurface states~\cite{liu2026physics}, and learned genetic operators have been obtained by meta-black-box optimisation~\cite{lange2023discovering}. On the local-search side, one- and $k$-opt flip search for binary quadratic problems~\cite{merz2002bqp}, iterated local search~\cite{lourenco2010ils}, and memetic and path-relinking methods with elite pools~\cite{merz2004memetic,glover2010ddts,lu2010hybrid,wang2012pathrelinking} are established; their flip-gain formulas rely on a quadratic objective, which the coupled max--min rate is not. For RIS, Zhang et al.~\cite{zhang2022configuring} gave an optimal method for the single-user uncoupled case, and Hammami et al.~\cite{hammami2026statistical} found single-flip coordinate descent the best trade-off on statistical multiport models.

\textbf{Coupled models and efficient updates.} Mutual impedance models~\cite{qian2021mutual}, multiport network frameworks~\cite{nerini2024universal,nossek2024physically}, S-parameter designs with full-wave validation~\cite{abrardo2024design} and closed-form optima under coupling~\cite{nerini2026global} establish the impedance-domain channel used here. PhysFad~\cite{faqiri2023physfad} and Prod'homme and del Hougne~\cite{prodhomme2023efficient} model RIS-parametrised environments with coupled dipoles, and the latter use low-rank updates to compute channel realisations efficiently. They derive the rank-$m$ update for configuration changes; we use the same algebra inside a deadline-bounded search that scans the complete pair neighbourhood at each step. Measured datasets~\cite{tewes2023comprehensive,rossanese2024open,piana2026brisc} and simulators~\cite{burtakov2023qris} supply channels but do not fix a deadline-bounded protocol for comparing configuration methods.

%======================================================================
\section{Problem and Coupled Evaluator}\label{sec:problem}
%======================================================================
\textbf{Channel.} The surface has $N$ thin-wire elements; each is terminated by one of two passive loads $z_0,z_1$ (1\,$\Omega$ resistance, reactances $\pm70\,\Omega$ around the isolated-wire resonance). Let $Z_{\mathrm{RIS}}$ be the full $N\times N$ impedance matrix from NEC-2~\cite{burke1981nec,pynec}, after a Schur-complement reduction over the 50\,$\Omega$-terminated BS and user ports. With $Z_L(\mathbf b)=\mathrm{diag}(z_{b_1},\dots,z_{b_N})$ and $T(\mathbf b)=(Z_{\mathrm{RIS}}+Z_L(\mathbf b))^{-1}$, the channel from the four BS ports to the $K$ users is
\begin{equation}
H(\mathbf b)=\Hd + C_{\mathrm{out}}\,T(\mathbf b)\,C_{\mathrm{in}},
\label{eq:operator}
\end{equation}
where $\Hd$ is the direct channel and $C_{\mathrm{in}}$, $C_{\mathrm{out}}$ couple the elements to the BS and the users. The off-diagonal entries of $Z_{\mathrm{RIS}}$ are the mutual coupling, which makes the effect of one element depend on all others.

\textbf{Objective.} A regularised zero-forcing precoder $W$ is computed once from the all-OFF channel and kept fixed. With effective field $U(\mathbf b)=H(\mathbf b)W$ and noise power $\sigma^2$,
\begin{equation}
\max_{\mathbf b\in\{0,1\}^N} f(\mathbf b),\qquad f(\mathbf b)=\min_{k}\log_2\Bigl(1+\frac{|U_{kk}|^2}{\sum_{j\ne k}|U_{kj}|^2+\sigma^2}\Bigr).
\label{eq:objective}
\end{equation}

\textbf{Exact single and pair moves.} Flipping element $i$ changes one diagonal load by $\delta_i=z_{1-b_i}-z_{b_i}$. Given $T$, precompute $L=C_{\mathrm{out}}T\in\mathbb C^{K\times N}$, $R=TC_{\mathrm{in}}W\in\mathbb C^{N\times S}$ and $d_i=1/\delta_i+T_{ii}$. By Sherman--Morrison, the field after flipping $i$ is
\begin{equation}
U^{(i)}=U-\frac{\ell_i\,r_i^{\top}}{d_i},
\end{equation}
with $\ell_i$ the $i$-th column of $L$ and $r_i^\top$ the $i$-th row of $R$. For a pair $(i,j)$, the $2\times2$ Woodbury identity gives, with $\Delta_{ij}=d_id_j-T_{ij}T_{ji}$,
\begin{equation}
U^{(ij)}=U-\ell_i u_i^{\top}-\ell_j u_j^{\top},\quad u_i=\frac{d_jr_i-T_{ij}r_j}{\Delta_{ij}},\quad u_j=\frac{d_ir_j-T_{ji}r_i}{\Delta_{ij}}.
\label{eq:pair}
\end{equation}
\textbf{Cost.} Forming $L$ and $R$ costs $O(N^2(K+S))$ once per state. Each single move then costs $O(KS)$ and each pair move $O(KS)$ plus a $2\times2$ determinant, so a full scan of all $N(N-1)/2$ pairs costs $O(N^2KS)$. Scoring each pair by a fresh inverse would cost $O(N^3)$ per candidate, i.e.\ $O(N^5)$ per scan. Committing an accepted move updates $T$ by a rank-one or rank-two correction in $O(N^2)$. All methods in this paper, including the evolutionary baselines, use this evaluator. In our implementations the population methods score each individual by a full state construction: offspring usually differ from their parents in many bits, which limits what a low-rank update can save.

%======================================================================
\section{Method}\label{sec:method}
%======================================================================
\textbf{Exact pair search.} The base method, \textsc{pair} (Algorithm~\ref{alg:pair}), is an anytime 2-opt local search. From the current state it takes the best improving single flip; when no single flip improves, it scans all pairs with~\eqref{eq:pair} and takes the best improving pair; when neither improves, the state is a 2-opt local optimum and the search restarts. The incumbent is the best state seen. The search starts from all OFF, so its result is never worse than the reference.

\begin{algorithm}[t]
\caption{Deadline-bounded exact pair search with a restart rule (\textsc{pair}, b3, b5, \bts{})}\label{alg:pair}
\footnotesize
\begin{algorithmic}[1]
\Require evaluator $\mathcal E$ for~\eqref{eq:operator}--\eqref{eq:objective}; absolute deadline $t_D$; seed; restart rule \textsc{Restart}
\State $s\gets\mathcal E.\textsc{State}(\mathbf 0)$; $s^\star\gets s$; $\mathcal P\gets\emptyset$ \Comment{all OFF; elite pool, at most 8}
\State $s\gets\textsc{Climb}(s)$; $r\gets0$
\While{$t_D-\textsc{Now}()>\tau_{\mathrm{restart}}$}
  \State $r\gets r+1$; $s\gets\mathcal E.\textsc{State}(\textsc{Restart}(r,\mathcal P))$ \Comment{one full $O(N^3)$ construction}
  \State $s\gets\textsc{Climb}(s)$
\EndWhile
\State \Return $s^\star$
\Statex
\Function{Climb}{$s$}
  \Loop
    \State $g\gets\mathcal E.\textsc{Singles}(s)$ \textbf{if} the scan can finish before $t_D$ \textbf{else} \textbf{break} \Comment{$N$ rank-one scores}
    \If{$\max g>f(s)$} $s\gets\mathcal E.\textsc{Flip}(s,\arg\max g)$; \textbf{continue} \EndIf
    \State $q\gets\mathcal E.\textsc{Pairs}(s)$ \textbf{if} the scan can finish before $t_D$ \textbf{else} \textbf{break} \Comment{$N(N{-}1)/2$ rank-two scores}
    \If{$\max q>f(s)$} $s\gets\mathcal E.\textsc{Flip}(s,\arg\max q)$; \textbf{continue} \EndIf
    \State \textbf{break} \Comment{2-opt local optimum}
  \EndLoop
  \State \textsc{Remember}$(\mathcal P,s)$; update $s^\star$ if $f(s)>f(s^\star)$; \Return $s$
\EndFunction
\Statex
\Function{Restart}{$r,\mathcal P$} \Comment{the four variants differ only here}
  \If{\bts{} or b5, $|\mathcal P|\ge2$ and $r\bmod3=0$}
    \State draw distinct $e,e'\in\mathcal P$; \Return $\mathbf c$ with $c_n=e_n$ w.p.\ $\tfrac12$, else $e'_n$ \Comment{elite recombination}
  \ElsIf{\bts{} or b3}
    \State $\rho\gets 18/64$ if $r$ odd, else $\rho\gets 1/2$; \Return $\mathbf c$ with $c_n\sim\mathrm{Bernoulli}(\rho)$ \Comment{restart density}
  \Else\ \Return $\mathbf c$ with $c_n\sim\mathrm{Bernoulli}(1/2)$ \Comment{\textsc{pair}, uniform}
  \EndIf
\EndFunction
\end{algorithmic}
\end{algorithm}

\textbf{Restart density (b3).} Uniform restarts put about half of the elements ON, whereas good states of this problem need not do so. b3 alternates two ON densities fixed before any evaluation, $18/64$ and $1/2$, so that half of the restarts explore sparser states.

\textbf{Elite recombination (b5).} The local optima reached during the climbs (single-flip optima after a pair move, and every 2-opt optimum) are offered to a pool $\mathcal P$ that keeps at most eight distinct states, ordered by value. Every third restart, once two elites exist, the start state is a uniform crossover of two distinct elites drawn at random. The child keeps the bits on which the elites agree and mixes the rest. Local search then repairs it, which places the restart near the good basins already found rather than in a random region.

\textbf{\bts{}.} \bts{} is a memetic iterated local search~\cite{lourenco2010ils,merz2004memetic}; recombination is its only evolutionary element. It combines the two rules with a fixed precedence: every third restart is a recombination (when $|\mathcal P|\ge2$), and the others alternate the two densities. Neither rule learns from the outcome of a run, and the search limits of \textsc{pair}, b3, b5 and \bts{} (the two densities, the pool size, the recombination period and the time reserves) were fixed in advance in a frozen configuration and were not tuned on the evaluation data.

\textbf{Deadline handling.} The deadline is absolute and checked with a monotonic clock. A scan is started only if its measured cost fits before the deadline; a scan that finishes late is discarded and recorded. Fixed time reserves keep room for the final state construction. Evaluator set-up and every scan are inside the timed section, so the reported latency is the complete controller time.

%======================================================================
\section{Baselines, Metrics and Protocol}\label{sec:protocol}
%======================================================================
\textbf{Baselines.} \textsc{greedy} is single-flip best-improvement search with uniform restarts, i.e.\ Algorithm~\ref{alg:pair} without the pair scan. \textsc{pair} is the base method, and b3 and b5 are the two ablations of \bts{}.

\emph{ME-GA} is our reimplementation of~\cite{zardi2024memory} from its published specification, with its population of \res{megaPop} individuals and its adaptive rates, adapted to 1-bit elements and a fixed precoder; it runs under the same deadline and returns its best individual so far. Because the paper publishes no code and leaves several operators unspecified, the reimplementation resolved \res{megaDecisions} documented gaps before any run, none tuned on results. The main ones are: (a)~the text's formula for the effectiveness index $\Theta$ uses a running minimum, which makes $\Theta\le0$ always and disables the memory-learning branch, so we follow the prose and use the per-iteration best; (b)~the memory-restoring probability of the published equation can leave $[0,\kappa_{\max}]$, so we normalise it by the largest $|\Theta|$ seen; (c)~selection, crossover and mutation are not defined, so we use binary tournament, one-point crossover and gene-wise bit flips; (d)~the original uses 3-bit elements and a per-step zero-forcing precoder. \res{megaSensItems} of these decisions are flagged for sensitivity analysis.

\emph{NSGA-II}~\cite{deb2002fast} treats the $K$ per-user rates as objectives and returns the individual with the largest worst-user rate seen in the call. It uses fast non-dominated sorting, crowding distance, binary crowded tournament, elitist $(\mu+\lambda)$ selection, uniform crossover with probability $p_c$ and bit-flip mutation at rate $1/N$; the initial population is the all-OFF state plus random individuals. Population size and $p_c$ were chosen from a grid of \res{nsgaGridN} settings ($P\in\{\res{nsgaGridP}\}$, $p_c\in\{\res{nsgaGridPc}\}$) by a rule declared before tuning, on separate development trajectories not used in this evaluation: $P=\res{nsgaPop}$, $p_c=\res{nsgaPc}$.

Both evolutionary baselines use the same coupled evaluator in batched form: children within \res{megaWoodM} flips of a base individual are scored by a rank-$m$ Woodbury update when at least \res{megaWoodMin} qualify, the others by linear solves. Table~\ref{tab:cost} summarises the work per step.

\begin{table}[t]
\centering
\caption{Work per step on the coupled evaluator ($N$ elements, $K$ users, $S$ streams, $m$ differing bits).}\label{tab:cost}
\scriptsize\setlength{\tabcolsep}{3pt}
\begin{tabular}{@{}llll@{}}
\toprule
Operation & Used by & Candidates & Cost\\
\midrule
state construction (solve) & restarts of all local searches; GA individuals & 1 & $O(N^3)$\\
rank-$m$ update from a base state & GA individuals within $m\le\res{megaWoodM}$ flips & 1 & $O(N^2m)$\\
precompute $L$, $R$ for a state & greedy, pair, b3, b5, \bts{} & -- & $O(N^2(K+S))$\\
single-flip scan, rank one & greedy, pair, b3, b5, \bts{} & $N$ & $O(NKS)$\\
pair scan, rank two & pair, b3, b5, \bts{} & $N(N{-}1)/2$ & $O(N^2KS)$\\
commit a move & greedy, pair, b3, b5, \bts{} & -- & $O(N^2)$\\
\bottomrule
\end{tabular}
\end{table}

For $N=64$, 128 and 256 the pair scan covers \res{nPairsSixtyFour}, \res{nPairsOneTwentyEight} and \res{nPairsTwoFiftySix} candidates, each scored at the cost of a few complex multiply--adds once $L$ and $R$ are known.

\subsection{Evaluation metrics}\label{sec:metrics}
Every run is described by three quantities.
\begin{itemize}
\item[\textbf{M1}] \emph{Solution quality}: the worst-user spectral efficiency $F(\mathbf b)=\min_k R_k(\mathbf b)$ of the returned state, in bit/s/Hz, where $R_k$ is user $k$'s rate under the fixed regularised zero-forcing precoder of the all-OFF channel, i.e.\ $f$ in~\eqref{eq:objective}.
\item[\textbf{M2}] \emph{Optimisation latency}: the measured wall-clock time of the search call in ms (\texttt{search\_elapsed\_ms}), including evaluator construction and any feature extraction or inference; we report its mean and 95th percentile.
\item[\textbf{M3}] \emph{Computational budget}: the hard deadline, at 5, 10, 50 and 100\,ms with 50\,ms primary, together with the number of candidate states the method scored within it.
\end{itemize}
The headline results relate M1 to M2 at each M3 level. Miss rates, bootstrap intervals and gains over all OFF are secondary.

\textbf{Protocol.} Each method runs as one Python process on a pinned CPU core with single-threaded numerical libraries. A run whose latency exceeds its deadline is a miss; a method that misses more than \res{vMaxMiss}\% of its runs at a budget is excluded from that budget's contrasts. The unit of analysis is the trajectory: frames and \res{vSeeds} seeds are averaged within a trajectory, methods are paired on the same trajectories, frames and seeds, and intervals come from a paired trajectory bootstrap with \res{vBoot} resamples, Bonferroni-corrected over a declared family of \res{vFamily} contrasts per budget. Data, methods, budgets, contrasts and decision rules were frozen in a dated declaration before the single pass over the sealed test trajectories. The contrasts of this paper are, at 50\,ms and separately for each array size:
\begin{itemize}
\item[H1] \bts{} $-$ \textsc{pair}: do restart density and elite recombination improve exact pair search?
\item[H2] \bts{} $-$ ME-GA: is the local search better than the reimplemented memory-enhanced GA at equal deadline?
\item[H5] \bts{} $-$ NSGA-II: is it better than the standard elitist evolutionary algorithm at equal deadline?
\end{itemize}
Each passes only if its Bonferroni interval lies entirely above zero and neither method is excluded by the miss rule. The ablation contrasts \bts{} $-$ b3, \bts{} $-$ b5 and \textsc{pair} $-$ \textsc{greedy} are descriptive.

%======================================================================
\section{Evaluation on GeoRIS-Bench}\label{sec:results}
%======================================================================
\textbf{Dataset.} We use the RIS-relevant scope (version 3) of GeoRIS-Bench~\cite{georisbench}: Sionna RT~\cite{hoydis2023sionna} paths over LoD2 buildings and terrain of \res{vMap}~\cite{peters2022automated,ledoux2019cityjson,ahn4}, joined through~\eqref{eq:operator} to NEC-2 models~\cite{burke1981nec,pynec} of $8\times8$, $8\times16$ and $16\times16$ thin-wire surfaces on building facades. Each RIS is placed \res{vDmin}--\res{vDmax}\,m from its BS, and every served user is blocked from the BS, visible from the RIS, within \res{vRadius}\,m of it and within \res{vAngle}$^\circ$ of the panel normal. Table~\ref{tab:setup} lists the setup; sites and users are drawn from geometry only, never from rates.

\begin{table}[t]
\centering
\caption{Benchmark setup (GeoRIS-Bench, RIS-relevant scope v3).}\label{tab:setup}
\scriptsize\setlength{\tabcolsep}{4pt}
\begin{tabular}{@{}ll@{}}
\toprule
City / sites & \res{vMap}; \res{vSites} RIS sites; BS--RIS distance \res{vDOneMin}--\res{vDOneMax}\,m (median \res{vDOneMed}\,m)\\
Trajectories & \res{vTraj} $\times$ \res{vFrames} frames (\res{vIncluded} with channels); split \res{vTrainTraj}/\res{vValTraj}/\res{vTestTraj} train/validation/sealed test\\
Users & 2--4 served per frame, \res{vServedUF} served user-frames; RIS--user distance \res{vDTwoLo}/\res{vDTwoMed}/\res{vDTwoHi}\,m (q10/median/q90)\\
Ray tracing & Sionna RT, \res{freqGHz}\,GHz, depth \res{vDepth}, LOS + specular + diffraction, $\res{vSamples}$ samples per source; one shared direct trace per frame\\
Surface & NEC-2 thin wires (24\,mm, 28/30\,mm pitch), $N=64/128/256$, two loads ($\pm70\,\Omega$), full $Z_{\mathrm{RIS}}$\\
Link & 4 BS ports, fixed RZF precoder of the all-OFF channel, 0.1\,W, 1\,MHz, 7\,dB noise figure\\
Headroom (validation) & median gain of a long search over all OFF: \res{vGainMed64} / \res{vGainMed128} / \res{vGainMed256}\,bit/s/Hz (64/128/256)\\
\bottomrule
\end{tabular}
\end{table}

\begin{table}[t]
\centering
\caption{Headline result, sealed test, primary 50\,ms budget: M1 worst-user SE (bit/s/Hz), M2 latency mean/p95 (ms), M3 candidate evaluations, per array size; contrasts with Bonferroni intervals and verdicts. Red P: pending.}\label{tab:headline}
\pendshorttrue\scriptsize\setlength{\tabcolsep}{2.6pt}
\begin{tabular}{@{}l|rrr|rrr|rrr@{}}
\toprule
 & \multicolumn{3}{c|}{64 elements} & \multicolumn{3}{c|}{128 elements} & \multicolumn{3}{c}{256 elements}\\
Method & M1 & M2 & M3 & M1 & M2 & M3 & M1 & M2 & M3\\
\midrule
greedy & \res{vRateGreedyFifty64} & \res{vLatGreedy64}/\res{vLatPGreedy64} & \res{vEvalsGreedy64} & \res{vRateGreedyFifty128} & \res{vLatGreedy128}/\res{vLatPGreedy128} & \res{vEvalsGreedy128} & \res{vRateGreedyFifty256} & \res{vLatGreedy256}/\res{vLatPGreedy256} & \res{vEvalsGreedy256}\\
pair & \res{vRatePairFifty64} & \res{vLatPair64}/\res{vLatPPair64} & \res{vEvalsPair64} & \res{vRatePairFifty128} & \res{vLatPair128}/\res{vLatPPair128} & \res{vEvalsPair128} & \res{vRatePairFifty256} & \res{vLatPair256}/\res{vLatPPair256} & \res{vEvalsPair256}\\
b3 & \res{vRateBthreeFifty64} & \res{vLatBthree64}/\res{vLatPBthree64} & \res{vEvalsBthree64} & \res{vRateBthreeFifty128} & \res{vLatBthree128}/\res{vLatPBthree128} & \res{vEvalsBthree128} & \res{vRateBthreeFifty256} & \res{vLatBthree256}/\res{vLatPBthree256} & \res{vEvalsBthree256}\\
b5 & \res{vRateBfiveFifty64} & \res{vLatBfive64}/\res{vLatPBfive64} & \res{vEvalsBfive64} & \res{vRateBfiveFifty128} & \res{vLatBfive128}/\res{vLatPBfive128} & \res{vEvalsBfive128} & \res{vRateBfiveFifty256} & \res{vLatBfive256}/\res{vLatPBfive256} & \res{vEvalsBfive256}\\
\bts{} (ours) & \res{vRateBtsFifty64} & \res{vLatBts64}/\res{vLatPBts64} & \res{vEvalsBts64} & \res{vRateBtsFifty128} & \res{vLatBts128}/\res{vLatPBts128} & \res{vEvalsBts128} & \res{vRateBtsFifty256} & \res{vLatBts256}/\res{vLatPBts256} & \res{vEvalsBts256}\\
ME-GA & \res{vRateMegaFifty64} & \res{vLatMega64}/\res{vLatPMega64} & \res{vEvalsMega64} & \res{vRateMegaFifty128} & \res{vLatMega128}/\res{vLatPMega128} & \res{vEvalsMega128} & \res{vRateMegaFifty256} & \res{vLatMega256}/\res{vLatPMega256} & \res{vEvalsMega256}\\
NSGA-II & \res{vRateNsgaFifty64} & \res{vLatNsga64}/\res{vLatPNsga64} & \res{vEvalsNsga64} & \res{vRateNsgaFifty128} & \res{vLatNsga128}/\res{vLatPNsga128} & \res{vEvalsNsga128} & \res{vRateNsgaFifty256} & \res{vLatNsga256}/\res{vLatPNsga256} & \res{vEvalsNsga256}\\
\midrule
H1 \bts{}$-$pair & \multicolumn{3}{c|}{\res{vHOne64} \res{vHOneVerdict64}} & \multicolumn{3}{c|}{\res{vHOne128} \res{vHOneVerdict128}} & \multicolumn{3}{c}{\res{vHOne256} \res{vHOneVerdict256}}\\
H2 \bts{}$-$ME-GA & \multicolumn{3}{c|}{\res{vHTwo64} \res{vHTwoVerdict64}} & \multicolumn{3}{c|}{\res{vHTwo128} \res{vHTwoVerdict128}} & \multicolumn{3}{c}{\res{vHTwo256} \res{vHTwoVerdict256}}\\
H5 \bts{}$-$NSGA-II & \multicolumn{3}{c|}{\res{vHFive64} \res{vHFiveVerdict64}} & \multicolumn{3}{c|}{\res{vHFive128} \res{vHFiveVerdict128}} & \multicolumn{3}{c}{\res{vHFive256} \res{vHFiveVerdict256}}\\
\bottomrule
\end{tabular}
\end{table}

\begin{table}[t]
\centering
\caption{M1 (worst-user SE, bit/s/Hz) at every budget level (M3), sealed test, and miss rates at 50\,ms. Red P: pending.}\label{tab:budgets}
\pendshorttrue\scriptsize\setlength{\tabcolsep}{2.2pt}
\begin{tabular}{@{}l|rrrr|rrrr|rrrr|rrr@{}}
\toprule
 & \multicolumn{4}{c|}{64: 5/10/50/100\,ms} & \multicolumn{4}{c|}{128: 5/10/50/100\,ms} & \multicolumn{4}{c|}{256: 5/10/50/100\,ms} & \multicolumn{3}{c}{miss \%, 50\,ms}\\
\midrule
greedy & \res{vRateGreedyFive64} & \res{vRateGreedyTen64} & \res{vRateGreedyFifty64} & \res{vRateGreedyHundred64} & \res{vRateGreedyFive128} & \res{vRateGreedyTen128} & \res{vRateGreedyFifty128} & \res{vRateGreedyHundred128} & \res{vRateGreedyFive256} & \res{vRateGreedyTen256} & \res{vRateGreedyFifty256} & \res{vRateGreedyHundred256} & \res{vMissGreedy64} & \res{vMissGreedy128} & \res{vMissGreedy256}\\
pair & \res{vRatePairFive64} & \res{vRatePairTen64} & \res{vRatePairFifty64} & \res{vRatePairHundred64} & \res{vRatePairFive128} & \res{vRatePairTen128} & \res{vRatePairFifty128} & \res{vRatePairHundred128} & \res{vRatePairFive256} & \res{vRatePairTen256} & \res{vRatePairFifty256} & \res{vRatePairHundred256} & \res{vMissPair64} & \res{vMissPair128} & \res{vMissPair256}\\
b3 & \res{vRateBthreeFive64} & \res{vRateBthreeTen64} & \res{vRateBthreeFifty64} & \res{vRateBthreeHundred64} & \res{vRateBthreeFive128} & \res{vRateBthreeTen128} & \res{vRateBthreeFifty128} & \res{vRateBthreeHundred128} & \res{vRateBthreeFive256} & \res{vRateBthreeTen256} & \res{vRateBthreeFifty256} & \res{vRateBthreeHundred256} & \res{vMissBthree64} & \res{vMissBthree128} & \res{vMissBthree256}\\
b5 & \res{vRateBfiveFive64} & \res{vRateBfiveTen64} & \res{vRateBfiveFifty64} & \res{vRateBfiveHundred64} & \res{vRateBfiveFive128} & \res{vRateBfiveTen128} & \res{vRateBfiveFifty128} & \res{vRateBfiveHundred128} & \res{vRateBfiveFive256} & \res{vRateBfiveTen256} & \res{vRateBfiveFifty256} & \res{vRateBfiveHundred256} & \res{vMissBfive64} & \res{vMissBfive128} & \res{vMissBfive256}\\
\bts{} (ours) & \res{vRateBtsFive64} & \res{vRateBtsTen64} & \res{vRateBtsFifty64} & \res{vRateBtsHundred64} & \res{vRateBtsFive128} & \res{vRateBtsTen128} & \res{vRateBtsFifty128} & \res{vRateBtsHundred128} & \res{vRateBtsFive256} & \res{vRateBtsTen256} & \res{vRateBtsFifty256} & \res{vRateBtsHundred256} & \res{vMissBts64} & \res{vMissBts128} & \res{vMissBts256}\\
ME-GA & \res{vRateMegaFive64} & \res{vRateMegaTen64} & \res{vRateMegaFifty64} & \res{vRateMegaHundred64} & \res{vRateMegaFive128} & \res{vRateMegaTen128} & \res{vRateMegaFifty128} & \res{vRateMegaHundred128} & \res{vRateMegaFive256} & \res{vRateMegaTen256} & \res{vRateMegaFifty256} & \res{vRateMegaHundred256} & \res{vMissMega64} & \res{vMissMega128} & \res{vMissMega256}\\
NSGA-II & \res{vRateNsgaFive64} & \res{vRateNsgaTen64} & \res{vRateNsgaFifty64} & \res{vRateNsgaHundred64} & \res{vRateNsgaFive128} & \res{vRateNsgaTen128} & \res{vRateNsgaFifty128} & \res{vRateNsgaHundred128} & \res{vRateNsgaFive256} & \res{vRateNsgaTen256} & \res{vRateNsgaFifty256} & \res{vRateNsgaHundred256} & \res{vMissNsga64} & \res{vMissNsga128} & \res{vMissNsga256}\\
\bottomrule
\end{tabular}
\end{table}

\begin{figure}[t]
\centering
\ifbench
\includegraphics[width=\textwidth]{figures/headline.pdf}
\else
\fbox{\parbox[c][2.4cm][c]{0.9\textwidth}{\centering\small M1 worst-user SE against M2 mean latency, one curve per method,\\ markers at the 5/10/50/100\,ms budgets (M3), one panel per array size.\\[2pt] \PENDING{headline figure, drawn by make\_results.py --bench}}}
\fi
\caption{Solution quality against latency at every budget, sealed test.}\label{fig:headline}
\end{figure}

\textbf{Results.} \PENDING{results text after the sealed test}. Tables~\ref{tab:headline} and~\ref{tab:budgets} and Fig.~\ref{fig:headline} report the single confirmatory pass. Three questions organise the reading: whether exact pair moves and the two restart rules improve over single flips and plain pair search on ray-traced urban channels (H1 and the ablations); how the evolutionary baselines compare at equal latency (H2, H5 and the M3 columns); and whether the larger arrays, whose pair neighbourhood grows from \res{nPairsSixtyFour} to \res{nPairsTwoFiftySix} moves, still leave time for restarts within 50\,ms.

%======================================================================
\section{Threats to Validity}\label{sec:threats}
%======================================================================
\textbf{Scope of the evidence.} The evaluation covers one city (\res{vMap}), \res{vSites} RIS sites and simulation only; no measured hardware is involved. Timings are Python on one pinned CPU core of a lab workstation, not a surface controller; every method is timed on the same core under the same rules.

\textbf{Baselines.} ME-GA is our reimplementation from the published specification, adapted to 1-bit elements and a fixed precoder; the original paper reports no numbers comparable to ours, so H2 compares with the specification, not with published results, and a different reading of its open operators could perform differently. NSGA-II hyper-parameters were tuned on separate development data, whereas the local searches use limits fixed in advance from a frozen configuration and were not tuned on GeoRIS-Bench; the latter may be far from optimal for 128 and 256 elements.

\textbf{Physical model.} The surface is modelled as perfectly conducting thin wires without substrate, ground plane or switch parasitics, and the 256-element model uses a single mesh. The precoder is fixed at the all-OFF channel and full channel knowledge is assumed.

\textbf{Scene content.} GeoRIS-Bench scenes contain LoD2 buildings and terrain only; vegetation, vehicles and street furniture are not modelled. This follows common practice in ray-traced RIS-optimisation studies, which use building-and-ground city geometry such as DeepMIMO urban scenarios~\cite{peng2025risnet,alkhateeb2019deepmimo} and Sionna RT scenes~\cite{hoydis2023sionna}. Because every controller is evaluated on the identical channel of each frame and compared through paired contrasts~\cite{glasserman1992crn}, an omitted scatterer changes the instance all methods solve but gives no method privileged information. \bts{} solves each frame from the current channel only, so transient vehicle blockage~\cite{boban2011vehicles,3gpp38901} or wind-driven foliage motion~\cite{pelet2004wind} cannot bias it through stale memory, whereas a memory-based controller such as ME-GA~\cite{zardi2024memory} would, if anything, be favoured by the omission. The main limitation is that foliage attenuation at 5\,GHz~\cite{itu2021p833,alnuaimi1998vegetation} is not always common to all paths: trees in front of the RIS or across the BS--RIS segment attenuate only the controllable part of the channel. Absolute RIS gains should therefore be read as upper bounds, and method differences may shrink in foliated streets.

%======================================================================
\section{Conclusion}
%======================================================================
We presented \bts{}, a deadline-bounded memetic iterated local search for coupled 1-bit RIS control that scores every single and pair flip through rank-one and rank-two updates and diversifies with predeclared restart densities and elite recombination. Its cost per pair candidate is a few multiply--adds instead of a linear solve, which lets a local search examine the complete pair neighbourhood repeatedly within tens of milliseconds. We evaluated it on GeoRIS-Bench against exact pair search, its ablations, a reimplemented ME-GA and NSGA-II, on worst-user rate, latency and budget. \PENDING{one-sentence summary of the sealed-test outcome}.

\subsubsection*{Disclosure of Interests.} The authors have no competing interests to declare that are relevant to the content of this article.

\bibliographystyle{splncs04}
\bibliography{references}
\end{document}
